\subsection{Tensor product of two complexes}

Let (C, d) be a complex of right A-modules and (C', d') a complex of left A-modules.

Endow the k-module \(C \otimes_{A} C'\) with the grading such that

\[
(\mathrm{C} \otimes_ {\mathbf {A}} \mathrm{C} ^ {\prime}) _ {n} = \sum_ {p + q = n} (\mathrm{C} _ {p} \otimes \mathrm{C} _ {q} ^ {\prime})
\]

and let D denote the unique k-linear endomorphism of degree \((-1)\) of C \(\otimes_{A} C'\) such that

\[
\mathbf {D} (x \otimes x ^ {\prime}) = d x \otimes x ^ {\prime} + (- 1) ^ {p} x \otimes d ^ {\prime} x ^ {\prime}, \quad x \in \mathrm{C} _ {p}, y \in \mathrm{C} _ {q} ^ {\prime}, p, q \in \mathbf {Z}. \tag {1}
\]

One has \(D \circ D = 0\) since, with the notations of (1),

\[
\mathrm{D} ^ {2} (x \otimes x ^ {\prime}) = d d x \otimes x ^ {\prime} + (- 1) ^ {p - 1} d x \otimes d ^ {\prime} x ^ {\prime} + (- 1) ^ {p} d x \otimes d ^ {\prime} x ^ {\prime} - x \otimes d ^ {\prime} d ^ {\prime} x ^ {\prime}.
\]

The complex of \(k\) -modules \((\mathbf{C} \otimes_{\mathbf{A}} \mathbf{C}', \mathbf{D})\) is called the tensor product complex of the complexes \((\mathbf{C}, d)\) and \((\mathbf{C}', d')\) .

Remarks. --- 1) When \(\mathbf{C}'\) is reduced to \(\mathbf{C}_0' = \mathbf{M}\) , then \((\mathbf{C} \otimes_{\mathbf{A}} \mathbf{C}')_n = \mathbf{C}_n \otimes_{\mathbf{A}} \mathbf{M}\) and \(D = d \otimes 1_M\) ; for example \(\mathbf{C} \otimes_{\mathbf{A}} \mathbf{A}_s\) is naturally identified with \(C\) . Similarly, when \(C\) is reduced to \(\mathbf{C}_0 = P\) , then \((\mathbf{C} \otimes_{\mathbf{A}} \mathbf{C}')_n = P \otimes_{\mathbf{A}} \mathbf{C}_n'\) and \(D = 1_P \otimes d\) .

\begin{enumerate}
\def\labelenumi{\arabic{enumi})}
\setcounter{enumi}{1}
\tightlist
\item
  For every integer \(r\) , we have \((C \otimes_{A} C')(r) = C(r) \otimes_{A} C'\) , but \((C \otimes_{A} C')(r)\) and \(C \otimes_{A} C'(r)\) do not in general have the same differential.
\end{enumerate}

Let \(p, q\) be two integers, \(x \in Z_p(C)\) , \(x' \in Z_q(C')\) ; then the element \(x \otimes x'\) of \(C_p \otimes C_q'\) belongs to \(Z_{p+q}(C \otimes C')\) by (1); moreover, if \(y \in C_{p+1}\) , \(y' \in C_{q+1}'\) , we have

\[
(x + d y) \otimes (x ^ {\prime} + d ^ {\prime} y ^ {\prime}) = x \otimes x ^ {\prime} + \mathrm{D} (y \otimes x ^ {\prime} + (- 1) ^ {p} (x + d y) \otimes y ^ {\prime})
\]

by passing to quotients, we deduce a k-linear map, called canonical

\[
\gamma_ {p, q} (\mathrm{C}, \mathrm{C} ^ {\prime}): \mathrm{H} _ {p} (\mathrm{C}) \otimes_ {\mathrm{A}} \mathrm{H} _ {q} (\mathrm{C} ^ {\prime}) \rightarrow \mathrm{H} _ {p + q} (\mathrm{C} \otimes_ {\mathrm{A}} \mathrm{C} ^ {\prime});
\]

if we equip \(\mathbf{H}(\mathbf{C})\otimes_{\mathbf{A}}\mathbf{H}(\mathbf{C}^{\prime})\) with the grading such that

\[
\left(\mathrm{H} (\mathrm{C}) \otimes \mathrm{H} (\mathrm{C} ^ {\prime})\right) _ {n} = \sum_ {p + q = n} \mathrm{H} _ {p} (\mathrm{C}) \otimes_ {\mathrm{A}} \mathrm{H} _ {p} (\mathrm{C} ^ {\prime}),
\]

the \(\gamma_{p,q}\) define a graded \(k\) -linear map of degree 0

\[
\gamma (\mathbf {C}, \mathbf {C} ^ {\prime}): \mathrm{H} (\mathbf {C}) \otimes_ {\mathbf {A}} \mathrm{H} (\mathbf {C} ^ {\prime}) \rightarrow \mathrm{H} (\mathbf {C} \otimes_ {\mathbf {A}} \mathbf {C} ^ {\prime}).
\]

Proposition 6 of II, p.~59, is reformulated as follows:

PROPOSITION 1. --- If the complexes C and C' are zero on the right, then C ⊗\_A C' is zero on the right, and the canonical k-linear map

\[
\gamma_ {0, 0} (\mathbf {C}, \mathbf {C} ^ {\prime}): \mathrm{H} _ {0} (\mathbf {C}) \otimes_ {\mathbf {A}} \mathrm{H} _ {0} (\mathbf {C} ^ {\prime}) \rightarrow \mathrm{H} _ {0} (\mathbf {C} \otimes_ {\mathbf {A}} \mathbf {C} ^ {\prime})
\]

is bijective.

Let \(u: (\mathbf{C}, d) \to (\mathbf{C}_1, d_1)\) be a morphism of complexes of right A-modules and \(u': (\mathbf{C}', d') \to (\mathbf{C}_1', d_1')\) a morphism of complexes of left A-modules; then \(u \otimes u': \mathbf{C} \otimes_{\mathbf{A}} \mathbf{C}' \to \mathbf{C}_1 \otimes_{\mathbf{A}} \mathbf{C}_1'\) is a morphism of complexes of \(k\) -modules; indeed, it is graded of degree 0, and if one denotes by D and \(\mathbf{D}_1\) the differentials of \(\mathbf{C} \otimes \mathbf{C}'\) and \(\mathbf{C}_1 \otimes \mathbf{C}_1'\) we have, for \(p, q \in \mathbf{Z}, x \in \mathbf{C}_p, x' \in \mathbf{C}_q\) ,

\[
\begin{array}{c} (u \otimes u ^ {\prime}) (\mathrm{D} (x \otimes x ^ {\prime})) = u (d x) \otimes u ^ {\prime} (x ^ {\prime}) + (- 1) ^ {p} u (x) \otimes u ^ {\prime} (d ^ {\prime} x ^ {\prime}) = \\ = d _ {1} u (x) \otimes u ^ {\prime} (x ^ {\prime}) + (- 1) ^ {p} u (x) \otimes d _ {1} ^ {\prime} u ^ {\prime} (x ^ {\prime}) = \mathrm{D} _ {1} (u (x) \otimes u ^ {\prime} (x ^ {\prime}))  . \end{array}
\]

Moreover, the following diagram is commutative:

\[
\begin{array}{c} \mathrm{H(C)} \otimes_ {\mathrm{A}} \mathrm {H(C^ {\prime})} \xrightarrow {\gamma (\mathrm{C,C} ^ {\prime})} \mathrm{H(C} \otimes_ {\mathrm{A}} \mathrm {C^ {\prime}}) \\ \mathrm{H(u)} \otimes \mathrm {H(u^{\prime})} \Big \downarrow \\ \mathrm {H(C_ {1})} \otimes_ {\mathrm{A}} \mathrm {H(C_ {1} ^ {\prime})} \xrightarrow [ \gamma (\mathrm {C_ {1}} , \mathrm {C_ {1} ^ {\prime}}) ]{} \mathrm {H(C_ {1}} \otimes_ {\mathrm{A}} \mathrm {C_ {1} ^ {\prime}}). \end{array}
\]

Let A° be the opposite k-algebra of A, and let C° (resp. C'°) be the complex C (resp. C') viewed as a complex of left (resp. right) A°-modules.

\[
\sigma (\mathbf {C}, \mathbf {C} ^ {\prime}): \mathbf {C} \otimes_ {\mathbf {A}} \mathbf {C} ^ {\prime} \rightarrow \mathbf {C} ^ {\prime \circ} \otimes_ {\mathbf {A} ^ {\circ}} \mathbf {C} ^ {\circ}
\]

the unique graded \(k\) -linear map of degree 0 such that, for \(x \in C_p\) , \(x' \in C_q'\) , \(p, q \in \mathbf{Z}\) , one has

\[
\sigma (\mathbf {C}, \mathbf {C} ^ {\prime}) (x \otimes x ^ {\prime}) = (- 1) ^ {p q} x ^ {\prime} \otimes x.
\]

PROPOSITION 2. — The map \(\sigma(\mathbf{C},\mathbf{C}'):\mathbf{C}\otimes_{\mathbf{A}}\mathbf{C}'\to\mathbf{C}^{\prime\circ}\otimes_{\mathbf{A}^{\circ}}\mathbf{C}^{\circ}\) is an isomorphism of complexes of \(k\) -modules, whose inverse isomorphism is \(\sigma(\mathbf{C}^{\prime\circ},\mathbf{C}^{\circ})\) .

Since the maps \(\sigma(\mathbf{C},\mathbf{C}^{\prime})\) and \(\sigma(\mathbf{C}^{\circ},\mathbf{C}^{\circ})\) are inverses of each other, it suffices to prove that \(\sigma(\mathbf{C},\mathbf{C}^{\prime})\) is a morphism of complexes. Now, for

\[
x \in \mathrm{C} _ {p} = \mathrm{C} _ {p} ^ {\circ}, \quad x ^ {\prime} \in \mathrm{C} _ {q} ^ {\prime} = \mathrm{C} _ {q} ^ {\prime \circ}, \quad p, q \in \mathbf {Z},
\]

we have, writing \(D^{c}\) for the differential of \(C' \otimes_{A^{0}} C\) ,

\[
\begin{array}{l} \sigma (\mathrm{C}, \mathrm{C} ^ {\prime}) \circ \mathrm{D} (x \otimes x ^ {\prime}) = \sigma (\mathrm{C}, \mathrm{C} ^ {\prime}) (d x \otimes x ^ {\prime} + (- 1) ^ {p} x \otimes d ^ {\prime} x ^ {\prime}) = \\ = (- 1) ^ {(p + 1) q} x ^ {\prime} \otimes d x + (- 1) ^ {p + p (q + 1)} d ^ {\prime} x ^ {\prime} \otimes x = (- 1) ^ {p q} d ^ {\prime} x ^ {\prime} \otimes x + (- 1) ^ {p q + q} x ^ {\prime} \otimes d x \\ = (- 1) ^ {p q} \mathrm{D} ^ {\circ} (x ^ {\prime} \otimes x) = \mathrm{D} ^ {\circ} \circ \sigma (\mathrm{C}, \mathrm{C} ^ {\prime}) (x \otimes x ^ {\prime}); \end{array}
\]

this gives \(\sigma(C, C') \circ D = D^{\circ} \circ \sigma(C, C')\) , whence the desired assertion.

The isomorphism \(\sigma(\mathbf{C},\mathbf{C}^{\prime}):\mathbf{C}\otimes_{\mathbf{A}}\mathbf{C}^{\prime}\to\mathbf{C}^{\circ}\otimes_{\mathbf{A}^{\circ}}\mathbf{C}^{-}\) is called the commutation isomorphism of the tensor product of the complexes \(\mathbf{C}\) and \(\mathbf{C}^{\prime}\) .

If \(u: \mathbf{C} \to \mathbf{C}_{1}\) and \(v: \mathbf{C}' \to \mathbf{C}_{1}'\) are two morphisms of complexes as above, we have a commutative diagram:

\[
\begin{array}{c} \mathrm{C} \otimes_ {\mathrm{A}} \mathrm{C} ^ {\prime} \xrightarrow {\sigma (\mathrm{C} , \mathrm{C} ^ {\prime})} \mathrm{C} ^ {\prime \circ} \otimes_ {\mathrm{A} ^ {\circ}} \mathrm{C} ^ {\circ} \\ u \otimes u ^ {\prime} \Bigg \downarrow \\ \mathrm{C} _ {1} \otimes_ {\mathrm{A}} \mathrm{C} _ {1} ^ {\prime} \xrightarrow [ \sigma (\mathrm{C} _ {1} , \mathrm{C} _ {1} ^ {\prime}) ]{} \mathrm{C} _ {1} ^ {\prime \circ} \otimes_ {\mathrm{A} ^ {\circ}} \mathrm{C} _ {1} ^ {\circ}. \end{array}
\]

Suppose for the remainder of this number that the ring A is commutative (cf.~no. 9 for the general case).

Let C, C′, C′′ be three complexes of A-modules; the canonical homomorphism of A-modules (III, p. 64)

\[
\varphi : \left(\mathrm{C} \otimes_ {\mathrm{A}} \mathrm{C} ^ {\prime}\right) \otimes_ {\mathrm{A}} \mathrm{C} ^ {\prime \prime} \rightarrow \mathrm{C} \otimes_ {\mathrm{A}} \left(\mathrm{C} ^ {\prime} \otimes_ {\mathrm{A}} \mathrm{C} ^ {\prime \prime}\right)
\]

is an isomorphism of complexes, as one verifies immediately using the definitions.

More generally, let \((\mathbf{C}^{(i)}, d^{(i)})_{i \in I}\) be a family of complexes of A-modules, where the set I is finite and totally ordered; for simplicity of notation, we will identify I with the interval \([1, r]\) of N. Endow the A-module \(C = \bigotimes_{i=1}^{r} C^{(i)}\) with the grading such that

\[
\mathrm{C} _ {n} = \sum_ {p _ {1} + p _ {2} + \dots + p _ {r} = n} (\mathrm{C} ^ {(1)}) _ {p _ {1}} \otimes (\mathrm{C} ^ {(2)}) _ {p _ {2}} \otimes \dots \otimes (\mathrm{C} ^ {(r)}) _ {p _ {r}},
\]

and let us define a graded A-endomorphism of degree (-1) of C by

\[
\mathrm{D} \left(x _ {1} \otimes \dots \otimes x _ {r}\right) = \sum_ {j = 1} ^ {r} (- 1) ^ {p _ {1} + \dots + p _ {j - 1}} x _ {1} \otimes \dots \otimes x _ {j - 1} \otimes d _ {j} x _ {j} \otimes x _ {j + 1} \otimes \dots \otimes x _ {r}
\]

where \(x_{i}\in (\mathbf{C}^{(i)})_{p_{i}}\) for \(i = 1,\dots,n\) . Then (C, D) is a complex of A-modules called the tensor product complex of the family \((\mathbf{C}_{i}, d_{i})\) . For every strictly increasing sequence \(r_{0}, \ldots, r_{k}\) of [0, r] such that \(r_{0} = 0\) , \(r_{k} = r\) , the canonical associativity isomorphism

\[
\bigotimes_ {j = 0} ^ {k - 1} \left(\bigotimes_ {i = r _ {j} + 1} ^ {r _ {j + 1}} \mathbf {C} ^ {(i)}\right)\rightarrow \bigotimes_ {i = 1} ^ {r} \mathbf {C} ^ {(i)}
\]

is an isomorphism of complexes.

We define as above a graded homomorphism of degree 0

\[
\gamma ((\mathbf {C} ^ {(i)})): \underset {i \in \mathbf {I}} {\otimes} \mathrm{H} (\mathbf {C} ^ {(i)}) \to \mathrm{H} \left(\underset {i \in \mathbf {I}} {\otimes} \mathbf {C} ^ {(i)}\right) .
\]

Remarks. --- 3) One can define the tensor product of a finite family of complexes without endowing the index set with a total order (X, p.~185, exercise 3).

\begin{enumerate}
\def\labelenumi{\arabic{enumi})}
\setcounter{enumi}{3}
\tightlist
\item
  Suppose that each \(\mathbf{C}^{(i)}\) is endowed with a graded algebra structure compatible with its grading and such that the \(d^{(i)}\) are antiderivations (III, p.~117). Let us then equip \(\bigotimes_{i\in I}\mathbf{C}^{(i)}\) with the structure of left graded tensor product algebra of the given structures (III, p.~49). Then D is an antiderivation. Indeed, using the associativity of the tensor product, one may assume that I = \{1, 2\}; let then \(p_1\) , \(q_1\) , \(p_2\) , \(q_2 \in \mathbf{Z}\) , \(x_1 \in (\mathbf{C}^{(1)})_{p_1}\) , \(y_1 \in (\mathbf{C}^{(1)})_{q_1}\) , \(x_2 \in (\mathbf{C}^{(2)})_{p_2}\) , \(y_2 \in (\mathbf{C}^{(2)})_{q_2}\) ; we have
\end{enumerate}

\[
\begin{array}{l} \left(\mathrm{D} (x _ {1} \otimes x _ {2})\right) (y _ {1} \otimes y _ {2}) + (- 1) ^ {p _ {1} + p _ {2}} (x _ {1} \otimes x _ {2}) (\mathrm{D} (y _ {1} \otimes y _ {2})) = \\ = (d x _ {1} \otimes x _ {2} + (- 1) ^ {p _ {1}} x _ {1} \otimes d x _ {2}) (y _ {1} \otimes y _ {2}) + \\ \quad + (- 1) ^ {p _ {1} + p _ {2}} (x _ {1} \otimes x _ {2}) (d y _ {1} \otimes y _ {2} + (- 1) ^ {q _ {1}} y _ {1} \otimes d y _ {2}) = \\ = (- 1) ^ {p _ {2} q _ {1}} (d x _ {1}) y _ {1} \otimes x _ {2} y _ {2} + (- 1) ^ {p _ {1} + (p _ {2} - 1) q _ {1}} x _ {1} y _ {1} \otimes (d x _ {2}) y _ {2} + \\ \quad + (- 1) ^ {p _ {1} + p _ {2} + p _ {2} (q _ {1} - 1)} x _ {1} d y _ {1} \otimes x _ {2} y _ {2} + (- 1) ^ {p _ {1} + p _ {2} + q _ {1} + p _ {2} q _ {1}} x _ {1} y _ {1} \otimes x _ {2} d y _ {2} \\ = (- 1) ^ {p _ {2} q _ {1}} [ (d x _ {1}) y _ {1} + (- 1) ^ {p _ {1}} x _ {1} d y _ {1} ] \otimes x _ {2} y _ {2} + \\ \quad + (- 1) ^ {p _ {1} + q _ {1} + p _ {2} q _ {1}} x _ {1} y _ {1} \otimes ((d x _ {2}) y _ {2} + (- 1) ^ {p _ {2}} x _ {2} d y _ {2}) \\ = (- 1) ^ {p _ {2} q _ {1}} [ d (x _ {1} y _ {1}) \otimes x _ {2} y _ {2} + (- 1) ^ {p _ {1} + q _ {1}} x _ {1} y _ {1} \otimes d (x _ {2} y _ {2}) ] \\ = (- 1) ^ {p _ {2} q _ {1}} \mathrm{D} (x _ {1} y _ {1} \otimes x _ {2} y _ {2}) = \mathrm{D} ((x _ {1} \otimes x _ {2}) (y _ {1} \otimes y _ {2}))  . \end{array}
\]

